\writeSectionFrame{\presentationTitle}{Das Problem}
\frameWithImageOnly{\BILDERPFAD/problem1.jpg}
\note{
	Man stelle sich eine Kreuzung vor. Hier gibt es jede Menge Akteure, die miteinander koordiniert
	werden müssen: Ampeln, Autos, Busse, Fußgänger, Bahnübergänge usw. 
}

\begin{frame}{Komplexe Kommunikation}
	\begin{center}
 		\includegraphics[scale=0.22]{\BILDERPFAD/problem2.png}
 	\end{center}
\end{frame}

\note{
	 An so einer Kreuzung entstehen sehr komplexe Kommunikationsvorgänge. Die Ampeln müssen
	 aufeinander abgestimmt sein. Die Autos müssen die Rot- und Grünphasen herausfinden können.
	 Die Ampeln wiederum können anhand des Verkehrsaufkommens gesteuert werden, müssen also
	 auch Informationen über die Autos bekommen können. Der Bahnübergang wiederum erfahren, 
	 wenn sich ein Zug nähert und dann die Autos oder sogar die Ampeln informieren.  
}


\begin{frame}{Komplexe Kommunikation}
	\begin{center}
 		\includegraphics[scale=0.22]{\BILDERPFAD/problem3.png}
 	\end{center}
\end{frame}

\note{
	 Die Lösung könnte eine Vermittler-Instanz sein, die alle Anfragen entgegennimmt und koordiniert. 
}